\section{Sử dụng phân bón hợp lí}

\subsection{Phản ứng của muối ammonium với base}

Muối ammonium phản ứng với base giải phóng khí ammonia. Đây là phản ứng quan trọng để giải thích vì sao không nên trộn trực tiếp phân ammonium với vôi hoặc các chất có tính base mạnh.

\begin{ghinho}{Phản ứng đặc trưng}
\begin{tasks}(2)
\task $\ce{NH4Cl + NaOH -> NaCl + NH3 ^ + H2O}$.
\task $\ce{NH4NO3 + NaOH -> NaNO3 + NH3 ^ + H2O}$.
\task $\ce{(NH4)2SO4 + 2NaOH -> Na2SO4 + 2NH3 ^ + 2H2O}$.
\task $\ce{2NH4Cl + Ca(OH)2 -> CaCl2 + 2NH3 ^ + 2H2O}$.
\end{tasks}
\end{ghinho}

Khí \ce{NH3} có mùi khai. Nếu nitrogen trong phân ammonium bị chuyển thành \ce{NH3} và thoát ra ngoài không khí thì lượng đạm cây nhận được giảm.

\begin{vidu}
Giải thích vì sao không nên trộn phân \ce{NH4Cl} với vôi tôi \ce{Ca(OH)2} rồi đem bón ngay.
\end{vidu}
\begin{loigiai}
Hai chất phản ứng:
\[
\ce{2NH4Cl + Ca(OH)2 -> CaCl2 + 2NH3 ^ + 2H2O}.
\]
Khí \ce{NH3} thoát ra làm thất thoát nitrogen của phân đạm.

\textbf{Vậy,} không nên trộn trực tiếp phân \ce{NH4Cl} với vôi tôi trước khi bón.
\end{loigiai}

\subsection{Phân bón và pH của đất}

Việc sử dụng lâu dài hoặc quá mức một số loại phân ammonium có thể làm pH đất giảm. Khi pH giảm, đất trở nên chua hơn.

\begin{notechem}
Cải tạo đất chua và bón phân là hai việc phải được thực hiện hợp lí. Không nên suy ra rằng có thể trộn trực tiếp mọi loại phân với vôi để trung hòa đất, vì một số phân ammonium sẽ phản ứng với base và làm thất thoát nitrogen dưới dạng \ce{NH3}.
\end{notechem}

\begin{vidu}
Một mẫu đất có pH giảm từ $6.2$ xuống $4.8$ sau nhiều năm bón phân không hợp lí. Nhận xét sự thay đổi môi trường đất.
\end{vidu}
\begin{loigiai}
Giá trị pH giảm và đều nhỏ hơn $7$, nên đất có tính acid mạnh hơn trước.

\textbf{Vậy,} đất đã trở nên chua hơn.
\end{loigiai}

\subsection{Phân bón và môi trường}

Bón quá nhiều phân không làm cây hấp thụ được toàn bộ chất dinh dưỡng. Phần dư có thể bị rửa trôi hoặc ngấm xuống đất và nguồn nước.

\begin{center}
\begin{tabularx}{0.96\textwidth}{Y Y}
\toprule
Tình huống & Hậu quả hoặc cách xử lí phù hợp \\
\midrule
Bón quá lượng & Tăng phần dinh dưỡng dư thừa, giảm hiệu quả sử dụng phân \\
Bón ngay trước mưa lớn & Phân dễ bị rửa trôi, gây thất thoát và ô nhiễm nguồn nước \\
Nước nhận nhiều hợp chất chứa N, P & Có thể thúc đẩy tảo phát triển mạnh, làm giảm chất lượng nước \\
Dùng phân ammonium lâu dài không hợp lí & Có thể làm pH đất giảm, đất chua hơn \\
Trộn phân ammonium với vôi & Có thể giải phóng \ce{NH3}, gây mất đạm \\
\bottomrule
\end{tabularx}
\end{center}

\subsection{Tính toán trong bài toán phân bón}

Bài toán phân bón không có cách giải riêng. Vẫn dùng bốn bước đã học:

\begin{ghinho}{Các bước giải}
\begin{enumerate}
\item Tính số mol khi bài toán có phản ứng hóa học; nếu chỉ dùng phần trăm khối lượng thì đổi trực tiếp về lượng chất dinh dưỡng cần thiết.
\item Viết phương trình hoặc sơ đồ phản ứng nếu có.
\item Tính các giá trị trung gian cần thiết: lượng chất dinh dưỡng, hiệu suất sử dụng, số mol khí, độ tinh khiết, ẩn số,...
\item Tính đại lượng đề hỏi và kết luận.
\end{enumerate}
\end{ghinho}

\begin{vidu}
Một vùng trồng trọt rộng \SI{12}{\hectare} cần thực tế \SI{80}{\kilo\gram} nitrogen được cây hấp thụ trên mỗi hectare. Dùng một loại urea chứa $96\%$ \ce{CO(NH2)2} theo khối lượng. Cây chỉ hấp thụ được $75\%$ lượng nitrogen bón vào. Phân được đóng bao \SI{50}{\kilo\gram}. Tính số bao tối thiểu cần mua.
\end{vidu}
\begin{loigiai}
Tổng khối lượng nitrogen cây cần hấp thụ:
\[
m_{N,\text{hấp thụ}}=12\times80=\SI{960}{\kilo\gram}.
\]
Tỉ lệ khối lượng nitrogen trong loại phân:
\[
w_N=0.96\times\frac{28}{60}=0.448.
\]
Khối lượng nitrogen cây hấp thụ được từ \SI{1}{\kilo\gram} phân:
\[
m_N=0.448\times0.75=\SI{0.336}{\kilo\gram}.
\]
Khối lượng phân cần dùng:
\[
m_{\text{phân}}=\frac{960}{0.336}\approx\SI{2857.14}{\kilo\gram}.
\]
Số bao:
\[
N=\frac{2857.14}{50}\approx57.14.
\]
Phải mua nguyên bao nên cần ít nhất $58$ bao.

\textbf{Vậy,} cần mua tối thiểu $58$ bao phân.
\end{loigiai}

\begin{tuluyen}
Viết phương trình hóa học và nêu hiện tượng khi đun nhẹ các chất sau với dung dịch \ce{NaOH} dư:
\begin{tasks}(2)
\task \ce{NH4Cl}.
\task \ce{NH4NO3}.
\task \ce{(NH4)2SO4}.
\task Một mẫu phân chứa ion \ce{NH4+}.
\end{tasks}
\dapso{Đều có khí \ce{NH3} mùi khai thoát ra; các PTHH tương ứng theo phản ứng của ion \ce{NH4+} với \ce{OH-}.}
\begin{loigiaiBT}
Ví dụ:
\[
\ce{NH4Cl + NaOH -> NaCl + NH3 ^ + H2O},
\]
\[
\ce{NH4NO3 + NaOH -> NaNO3 + NH3 ^ + H2O},
\]
\[
\ce{(NH4)2SO4 + 2NaOH -> Na2SO4 + 2NH3 ^ + 2H2O}.
\]
Trong cả ba trường hợp đều có khí ammonia mùi khai thoát ra.

\textbf{Vậy,} có thể dùng phản ứng với base để nhận biết ion ammonium trong nhiều mẫu phân đạm.
\end{loigiaiBT}
\end{tuluyen}
